\section{Introduction}

AI agent systems increasingly rely on coordination patterns such as planner-worker splits, reviewer
gates, specialist swarms, and distributed tool ownership. In practice, these topologies are still
chosen mostly heuristically. The central question of this paper is simple:
\emph{how should an agent architect choose among common coordination topologies when trading off
reliability, communication cost, and decomposition benefit?}

Our starting point is that many familiar coordination pathologies are architectural rather than purely
prompt-level failures. A reviewer can suppress some classes of error that a flat worker ensemble cannot.
A sequential task can become slower and less reliable when decomposed across multiple agents. A large
toolset can become harder to use when split across multiple specialists. These claims are familiar in
practice, but they are usually discussed informally and only after systems have already been built.

We use a minimal typed interface model inspired by wiring diagrams and interpret it epistemically:
topology determines what each worker or coordinator can observe, and those observation constraints induce
predictable costs and failure modes. Biology remains a motivating heuristic in the background, but it is
not the main claim of this paper. The motivating intuition is simply that local wiring patterns can
stabilize or destabilize distributed information-processing systems, a perspective that has proved useful
in the study of biological network motifs \cite{milo2002network,alon2007network}.

\paragraph{Contributions.}
This paper makes three contributions:
\begin{enumerate}[leftmargin=*]
\item It introduces a minimal typed and epistemic model for agent coordination, with enough structure to
reason about what different topologies make visible to workers and coordinators.
\item It derives predictive results for reviewer-gated execution, error amplification, sequential
handoffs, parallel decomposition, and tool-density scaling.
\item It compares those qualitative predictions to large-scale external evaluations of multi-agent
systems and maps them to an executable coordination substrate in the Operon implementation.
\end{enumerate}
